\chapter{Abstract}
\begin{onehalfspacing}
\begin{tabular}{@{} >{\bfseries}p{0.2\textwidth} p{0.77\textwidth} @{}}
\MakeUppercase{Keywords} &
Multi-echelon inventory optimisation; discrete-event simulation; digital twin;
transport consolidation; Bayesian optimisation; TPE; bullwhip effect; supply
chain; joint optimisation; Pareto frontier\\
\end{tabular}
\par
Classical multi-echelon inventory models optimise stock levels in isolation,
treating transport as a linear per-unit cost.  In practice, vehicle-based
transport has a step-function cost structure: a half-full truck costs nearly
as much to run as a full one.  When transport dominates supply-chain cost,
this mismatch makes inventory-only optimisation fundamentally suboptimal.

This project develops a modular, open-source \textbf{digital-twin simulation
framework} for Multi-Echelon Inventory Optimisation (MEIO) that jointly
optimises inventory and transport decisions.  The framework implements seven
inventory control policies, an arbitrary DAG network topology, and a transport
layer that enforces simultaneous volume and weight capacity constraints.  A
Tree-structured Parzen Estimator (TPE) optimiser searches a mixed-integer
space of stock levels and dispatch thresholds; a two-objective NSGA-II mode
maps the full cost--service Pareto frontier.

Applied to a 5-SKU, 1-supplier--2-warehouse--3-retailer network driven by
real M5 Walmart demand, joint optimisation achieves a 37\% total cost
reduction over the analytical Clark-Scarf baseline while simultaneously
raising fill rate above the 92\% service-level target --- results that
neither inventory-only nor transport-only optimisation can match.  In a
transport-cost-dominated network, extra warehouse safety stock enables
vehicle consolidation and generates substantial savings in transport
and shortage costs.

Extension experiments confirm the breadth of the result.  The elevated safety
stock incidentally absorbs a two-week supply disruption better than the
analytical baseline (E5), and policy comparison on equal transport footing
shows base-stock remains most cost-effective (E6).  The parameters are robust
to forecast error (E7), and the per-echelon bullwhip metric confirms
negligible amplification under base-stock with deterministic lead times (E8).
Stochastic testing (E9) confirms the parameters remain effective under
Poisson demand.  The framework is validated by 84 unit and integration tests
and four preliminary experiments (EV1--EV4), and is fully reproducible.
\par\end{onehalfspacing}
